\documentclass{beamer}
\input{tbook-deck}
\title{Permuted rotor loops without the star condition}
\author{Haobo Ma \and Wenlin Zhang}
\institute{ChronoAI Pte Ltd, Singapore \\ National University of Singapore (Wenlin Zhang)}
\date{Version 5 seminar materials}
\begin{document}
% Manuscript: title and §1
\begin{frame}
  \titlepage
  \note{The talk follows the accepted version by Haobo Ma and Wenlin Zhang. The question is Jenzer's Remark 5.5, p. 19 of arXiv:2610.21593v1. Distinguish the forward chain from the closed loop.}
\end{frame}
% Manuscript: §§1–3
\begin{frame}{Outline}
  \tableofcontents
  \note{Start with three planes. Then count selected pairs and explain why the only potentially nonzero off-diagonal group cancels. Close with the scope and a boundary example.}
\end{frame}
\section{The question}
\subsection{A non-star family}
% Manuscript: §1 and the triangle following Corollary 7
\begin{frame}{Three planes break the star condition}
  In $\mathrm{Cl}_{0,3}$, take $B_1=e_1e_2$, $B_2=e_1e_3$, $B_3=e_2e_3$.
  \[
    B_i^2=-1,\qquad B_iB_j=-B_jB_i\ (i\ne j),\qquad B_1B_2B_3=-1.
  \]
  \begin{itemize}
    \item The star condition would require every nonempty subset product to have zero scalar part.
    \item The forward chain can retain the triple term.
    \item Does the loop formula still hold for every return order?
  \end{itemize}
  \note{Since $B_1B_2=B_3$, the triple product is $-1$. The scalar of the forward chain is $c_1c_2c_3-s_1s_2s_3$. Jenzer's Theorem 5.2 handles the same-order loop; the extension concerns arbitrary return order without the star hypothesis.}
\end{frame}
\section{The scalar formula}
\subsection{Signs and a worked example}
% Manuscript: Theorem 1 and Corollary 9
\begin{frame}{Select a subset, then count pairs}
  \begin{bookresult}{loop-formula}
    For pairwise anticommuting square-$-1$ bivectors and any return permutation, select $i$ independently with weight $w_i=\sin^2(\theta_i/2)$.
    The loop scalar is the average of $(-1)^{e_\tau(S)}$, where $e_\tau(S)$ counts pairs in $S$ whose relative order is unchanged.
  \end{bookresult}
  \note{Use a finite-dimensional real Clifford algebra with nondegenerate quadratic form of any signature. The exact weight is $\prod_{i\in S}w_i\prod_{i\notin S}(1-w_i)$. The parity of $|S|(|S|+3)/2$ equals that of $\binom{|S|}{2}$; adding the inversion count has the same parity effect as subtracting it.}
\end{frame}
% Manuscript: Theorem 1 and Corollary 9, worked specialization
\begin{frame}{A return of $2,1,3$}
  \begin{itemize}
    \item Pairs keeping their order: $\{1,3\}$ and $\{2,3\}$.
    \item For $w_1=w_2=w_3=1/2$, subsets without $3$ all give $+1$.
    \item Subsets with $3$ give $+1,-1,-1,+1$: their contributions cancel.
  \end{itemize}
  Thus $t(\Theta_3(2,1,3))=1/2$. Reverse return instead gives $1$.
  \note{List the four choices of $1$ and $2$ conditional on selecting $3$. The whole return is not claimed to be the identity for order $2,1,3$; only the scalar is calculated. Reverse order cancels the factors themselves.}
\end{frame}
\section{Why the formula survives}
\subsection{Trace and exterior cancellation}
% Manuscript: Lemma 2 and Theorem 1 proof
\begin{frame}{The only off-diagonal exception}
  Expand by outward subset $S$ and return subset $T$.
  \begin{itemize}
    \item $S=T$: each repeated index squares; compute the rearrangement sign.
    \item $S\ne T$: cyclicity lets us conjugate by $B_j$ without changing the scalar.
    \item A sign flip forces zero except when $k$ is odd and $T=I\setminus S$.
  \end{itemize}
  \note{The residual set is $D=S\mathbin{\triangle}T$. If its positive size is even choose $j\in D$; if odd and proper choose $j\notin D$. Conjugation fixes $B_j$ and negates all other bivectors. In the exceptional group every coefficient contains $\prod_i c_i s_i$.}
\end{frame}
% Manuscript: Lemma 3 and complementary-subset correction
\begin{frame}{Swap with an even correction}
  Switch to exterior markers: $x^2=y^2=0$, $xy=-yx$.
  \[
    (1+x)(1+y)=(1+2xy)(1+y)(1+x).
  \]
  \begin{enumerate}
    \item Each swap contributes a commuting even factor.
    \item Reorder the return to reverse and cancel from the middle.
    \item The remaining product is even, so the full odd coefficient is zero.
  \end{enumerate}
  \note{The markers are a sign bookkeeping device, not new bivectors in the Clifford algebra. Their top coefficient sums the signs of all complementary choices; it multiplies $t(B_I)\prod_i c_i s_i$. The sum vanishes even when $t(B_I)$ does not.}
\end{frame}
\section{Scope and sources}
% Manuscript: Corollaries 4–9 and coverage section
\begin{frame}{What the scalar result does not say}
  \begin{itemize}
    \item The mixed-square and tracial formulas are further results, not claims about probability weights.
    \item The graph gives the sharp range and the factorization cuts in the negative-square case.
    \item Scalar powers for at most three factors follow from a quadratic identity; four factors can violate it.
  \end{itemize}
  \note{One four-factor loop has $t(R)=-1/4$ and $t(R^2)=-1/4$, not $T_2(-1/4)=-7/8$. The main frozen Lean result does not in itself certify the mixed-square, power, or graph extensions.}
\end{frame}
% Manuscript: introduction and bibliography
\begin{frame}{Sources}
  \begin{itemize}
    \item Haobo Ma and Wenlin Zhang, v5, \href{https://doi.org/10.5281/zenodo.24404942}{DOI 10.5281/zenodo.24404942}.
    \item N. Jenzer, arXiv:2610.21593v1, Theorems 5.2 and 5.5, Remark 5.5 (p. 19).
    \item Shirokov on conjugation averaging; Diehl--Reizenstein and de Bruijn on related exterior cancellation.
  \end{itemize}
  \note{The exact-version DOI refers to the paper, not this deck. The new contribution is the arbitrary-order cancellation without star, with further consequences as explained in v5.}
\end{frame}
\appendix
% Manuscript: Coverage and formal statements
\FormalVerification
\end{document}
